\documentclass[11pt]{article}
\usepackage[margin=1in]{geometry}
\usepackage{booktabs,array,enumitem,microtype}
\usepackage[hidelinks]{hyperref}
\usepackage{xcolor}
\setlength{\parskip}{4pt}\setlength{\parindent}{0pt}
\newcommand{\code}[1]{\texttt{#1}}
\newcolumntype{R}[1]{>{\raggedright\arraybackslash}p{#1}}
\title{Where TCSlib does not reach:\\ near misses and total misses in the STOC 2026 challenge corpus}
\author{lean\_formalization pipeline notes}
\date{5 November 2026}
\begin{document}
\maketitle

\section{Setup and headline numbers}

The corpus is the 259 Lean challenge files generated from the \code{tcs\_v3} STOC 2026 parquets
(199 proof-writing, 60 proof-strategy nodes). Each file is a self-contained statement
(definitions plus \code{sorry}-stubbed theorems) that compiles against the pinned TCSlib checkout
(Lean v4.25.0, Mathlib \code{029db123ddaa}). A keyword scan over twenty TCSlib concept areas
flagged files whose local definitions overlap a TCSlib definition; two passes of headless
refactor agents (66 agents, \$25) then rewrote every file where a TCSlib definition could replace a local one
without changing the statement. After that pass the corpus splits three ways:

\begin{center}
\begin{tabular}{lrrR{7.6cm}}
\toprule
Group & Files & Share & Meaning \\
\midrule
Uses TCSlib & 61 & 24\% & imports a TCSlib module and states at least one theorem through its definitions \\
Near miss & 37 & 14\% & defines an object TCSlib has, in a shape the paper cannot use \\
Total miss & 161 & 62\% & no local definition corresponds to anything in TCSlib \\
\midrule
 & 259 & 100\% & \\
\bottomrule
\end{tabular}
\end{center}

The scan originally flagged 61 near-miss files. Nine of them were plain renamings and were
converted in the second refactor pass; fifteen were keyword false positives (``walk'', ``game'',
``Gaussian'', ``restriction'' used in another sense) and are counted with the total misses here.
Earlier notes quoted 68 files (26\%) as using TCSlib; that figure double-counted seven files whose
agent declined to change anything. The count by actual imports is 61.

\section{Near misses: 37 files, fourteen encoding gaps}

A near miss is a file that reinvents a TCSlib object because the TCSlib version is indexed,
bundled or typed differently than the paper needs. Each gap below names the smallest TCSlib
change that would let the files use the library definition. Counts are entries in the gap report;
a few files appear under two gaps.

\begin{center}
\small
\begin{tabular}{@{}R{0.5cm}R{4.3cm}rR{5.3cm}R{4.3cm}@{}}
\toprule
 & Gap & Files & Why TCSlib's version does not fit & Fix \\
\midrule
C1 & Codes & 5 & \code{Codeword n $\alpha$} is \code{Fin n $\to$ $\alpha$} and the section needs \code{[Field $\alpha$]}; papers index words by cube points or an arbitrary finite type & metric definitions over \code{$\iota \to \alpha$} with \code{[Fintype $\iota$]}, plus \code{relDist} \\
C2 & CNF over generic variables; proof systems & 2 & \code{Literal n} is \code{Fin n}-indexed, clauses are lists; no resolution or Frege & \code{Literal (V : Type)}, finset clauses, a resolution derivation with width \\
C3 & Decision trees & 2 & \code{DecisionTree n} fixes inputs and leaves to \code{Bool} & parameters for input alphabet and leaf type; \code{RandDecisionTree} as a \code{FinDist} \\
C4 & Circuits & 7 & \code{AlgCircuit} has only the constant \code{one}; \code{DAGCircuit} has no constant wires and carries proof fields; no quantum circuits & \code{AlgGate.const}, a constant gate kind, an index-type parameter; quantum is a new area \\
C5 & Multigraphs & 4 & Karger's \code{Multigraph} is loopless with a multiset of edges & index-based multigraph \code{ends : Fin m $\to$ Sym2 V} with degree, volume, cut edges, conductance \\
C6 & Matchings & 4 & \code{Core.Matchings} lives on \code{SimpleGraph.Subgraph} & \code{IsMatching} on edge finsets, hyperedges and fractional matchings \\
C7 & Spanning trees and cut edge sets & 5 & cuts are capacities on a \code{Network}; no spanning-tree notion & \code{cutEdges G S} and one \code{IsSpanningTreeOf} on \code{Finset (Sym2 V)} \\
C8 & Property testers & 6 & BLR and low-degree tests exist, but no tester type & \code{Tester X Y} as a distribution over query trees, \code{acceptProb}, influence on product spaces \\
C9 & Expanders & 2 & expansion is spectral only & conductance, Cheeger, existence of explicit expanders \\
C10 & Protocols & 2 & two-party deterministic, public- and private-coin only & $k$-party number-in-hand protocols; biased-coin messages \\
C11 & Hermite functions; Reed--Solomon & 2 & probabilists' Hermite only; no Reed--Solomon or Reed--Muller codes & physicists' family with the rescaling lemma; both code families \\
C12 & Total variation & 2 & \code{statDist} is finite only & \code{tvDist} on \code{Measure} with the finite case as a lemma \\
C13 & Computed distributions & 5 & \code{FinDist} bundles its proofs, so a definition cannot build one & constructors for conditioning, kernel powers, $k$-fold products, uniform on a subset \\
C14 & Learning & 2 & external regret for finite experts, minimax, mistake bounds only & $\Phi$-regret; a PAC learner over measures \\
\bottomrule
\end{tabular}
\end{center}

Ranked by files unlocked per definition added, the order to work in is: index-generic codes (C1,
5 files, a few lines); spanning trees and cut edge sets (C7, 5); edge-finset and hypergraph
matchings (C6, 4); an index-based multigraph with loops (C5, 4, which also subsumes Karger's);
constant gates in both circuit models (C4 classical half, 4); \code{FinDist} constructors (C13, 5
files already half converted); a tester type (C8, 6). Together these fourteen items would move
about 35 of the 37 near-miss files into the ``uses TCSlib'' group. Each is a generalisation of
an existing definition, not new mathematics, and most are under a page of Lean.

\section{Total misses: 161 files, five areas}

A total miss defines nothing TCSlib defines. Grouping the 161 files by the paper's research area
(assigned by hand from the title, so the boundaries are approximate):

\begin{center}
\begin{tabular}{@{}R{6.2cm}rrR{6cm}@{}}
\toprule
Area & Files & Papers & What a library would need \\
\midrule
Continuous optimisation, convex geometry, lattices, numerical linear algebra, algebra and number theory & 50 & 21 & nothing TCS-specific: these statements live on Mathlib's analysis, linear algebra and lattice theory, and TCSlib is the wrong target \\
Learning, statistics, sampling, Markov chains, hashing, property testing, hardness of approximation & 46 & 24 & PAC and agnostic learners over measures, sample complexity, mixing times, measure-level total variation, hash families, CSP and PCP vocabulary \\
Graph algorithms: shortest paths, distance oracles, planar and minor-free graphs, treewidth, dynamic and streaming algorithms & 41 & 23 & weighted shortest-path distances and distance oracles, tree decompositions, planarity, a streaming and dynamic-algorithm model \\
Proof complexity, circuit lower bounds, pseudorandomness, cryptography, computability & 14 & 10 & resolution and Frege systems, algebraic complexity, explicit constructions, a Turing-machine model \\
Quantum & 10 & 5 & states, unitaries, Pauli and Clifford groups, measurements, tomography \\
\midrule
 & 161 & 83 & \\
\bottomrule
\end{tabular}
\end{center}

Two observations separate this group from the near misses.

\begin{itemize}[leftmargin=1.4em]
\item \textbf{New areas, not adjusted definitions.} A near miss closes with a one-line
generalisation. A total miss needs a vocabulary that does not exist in TCSlib, and for the 50
optimisation and geometry files the right home is Mathlib rather than TCSlib.
\item \textbf{Most content is paper-specific regardless.} Files average 14 to 16 local definitions
in every group (13.6 among total misses, 16.0 among near misses, 16.5 among files that already
use TCSlib). A library absorbs the shared vocabulary (a mixing time, a tree decomposition, a PAC
learner) but not the per-paper potential functions, gadgets and reductions that make up the bulk
of each statement. Even with the five areas above added, these files would look like the
current 24\%: a handful of library names at the top and a page of local definitions below.
\end{itemize}

If TCSlib were to grow toward this corpus, the order by files unlocked is a learning and
sampling area (46 files), a graph-algorithms area on weighted graphs (41), then quantum
(10 here plus 5 near-miss files, the largest single gap by file count but a new library rather
than a fix). The optimisation and geometry block should be left to Mathlib.

\section{Method and caveats}

Membership was computed on 5 November 2026 from the files in
\code{challenges/stoc2026\_pw} and \code{challenges/stoc2026\_ps}: a file ``uses TCSlib'' if it
has an \code{import TCSlib.*} line; the near-miss set is the concept scan's 61 flagged files
minus the nine converted and the fifteen false positives; everything else is a total miss. The
gap categories C1--C14 come from reading each near-miss file's definitions against the pinned
TCSlib source and are recorded with per-file detail in \code{notes/tcslib\_gap\_report.md}. The
area assignment of total misses is by paper title and would shift a few files under a different
rubric (for example, hardness-of-approximation papers sit with learning and statistics here, and
lattice-based quantum lower bounds sit with lattices). The per-file refactor notes are in
\code{batches/refactor\_tcslib/refactor\_notes\_all.jsonl} and are reproduced in each file's
provenance section in the public fork
(\url{https://github.com/qedsphere/tcslib}, branch \code{challenges/stoc2026}).

\end{document}
